\documentclass[11pt]{article}
\usepackage[margin=1in]{geometry}
\usepackage[T1]{fontenc}
\usepackage{lmodern}
\usepackage{amsmath,amssymb}
\usepackage{hyperref}

\title{Literature Status for arXiv:1808.00859: Pelczynski Space and GCCR}
\author{Run fa\_banach\_001, agent\_super\_00}
\date{2026-07-01}

\begin{document}
\maketitle

\section*{Status}

This is a \texttt{literature\_already\_answered} packet for a scoped
subquestion in Garcia-Lirola--Prochazka,
\emph{Pelczynski space is isomorphic to the Lipschitz free space over a
compact set}, arXiv:1808.00859.  The supporting paper is Hajek--Medina,
\emph{Compact retractions and Schauder decompositions in Banach spaces},
arXiv:2106.00331.

\section*{Source Question}

After proving that there is a compact convex set \(K\subset \mathbb P\) such
that the Pelczynski space \(\mathbb P\) is isomorphic to \(\mathcal F(K)\),
arXiv:1808.00859 observes that the proof could alternatively be completed by
showing that this \(K\) is a Lipschitz retract of \(\mathbb P\).  The authors
state that they do not know whether this is true, and then recall the broader
Godefroy--Ozawa question asking whether every separable Banach space admits a
compact convex generating Lipschitz retract.

\section*{Supporting Answer}

The later paper arXiv:2106.00331 studies exactly this compact-retraction
problem.  Its introduction states that every Banach space with a finite
dimensional decomposition admits a generating convex compact Lipschitz
retract.  It further states that if the space has a monotone Schauder basis,
then such a retract can be chosen with Lipschitz constant \(1+\varepsilon\)
for every \(\varepsilon>0\).

In Section ``Compact Lipschitz retracts'', the authors explicitly identify the
connection with arXiv:1808.00859: their FDD construction gives a positive
answer to the question asked in their citation \texttt{GP19}, namely whether
Pelczynski space has a GCCR.  Since Pelczynski space has a Schauder basis, it
falls under the FDD/Schauder-basis theorem.

\section*{Scope}

The answer is scoped.  It shows that Pelczynski space admits some generating
convex compact Lipschitz retract.  It does not prove that the particular
compact set \(K=\overline{\operatorname{conv}}\{e_n/n:n\in\mathbb N\}\) used
in arXiv:1808.00859 is itself a Lipschitz retract.  It also does not solve the
full Godefroy--Ozawa problem for arbitrary separable Banach spaces, nor the
separate compact-model question asking whether every separable
infinite-dimensional \(X\) admits a compact metric \(K\) with
\(\mathcal F(X)\simeq \mathcal F(K)\).

\section*{Search Evidence}

The deterministic lane queue selected arXiv:1808.00859.  Cheap run indexes
were searched for the arXiv id, title variants, Pelczynski-space keywords, and
Lipschitz-free compact-model phrases.  No exact packet for arXiv:1808.00859
was found.  A local full-source search for compact convex generating
Lipschitz retracts found arXiv:2106.00331, which explicitly cites the source
paper as \texttt{GP19} and states the Pelczynski-space GCCR consequence.
Existing run memory already contains a broader no-promotion attempt on the
general Godefroy--Ozawa problem, so this packet does not repackage that global
open problem as solved.

\begin{thebibliography}{9}
\bibitem{GP19}
L.~C. Garcia-Lirola and A.~Prochazka,
\emph{Pelczynski space is isomorphic to the Lipschitz free space over a
compact set}, arXiv:1808.00859.

\bibitem{HM21}
P.~Hajek and R.~Medina,
\emph{Compact retractions and Schauder decompositions in Banach spaces},
arXiv:2106.00331.
\end{thebibliography}

\end{document}
